% !TeX document-id = {6b928c0d-85af-41d1-a21b-e6c7a23a6e8d}
% !TeX TXS-program:compile = txs:///pdflatex/[--shell-escape]
% !TeX root = 00_Orga.tex
% <<<<<<< HEAD
% \documentclass[accentcolor=3c,landscape,ngerman,presentation,t,usenames,dvipsnames,svgnames,table,colorframetitle=true]{tudabeamer}
% =======

%include guard:
\ifdefined\GlobalAlreadyIncluded
    \expandafter\endinput
\fi

\gdef\GlobalAlreadyIncluded{}
%

\documentclass[accentcolor=3c,landscape,ngerman,presentation,t,usenames,dvipsnames,svgnames,table, aspectratio=169,colorframetitle=true]{tudabeamer}
%>>>>>>> online_rebase

% Template-Modifikationen
\addtobeamertemplate{frametitle}{}{\vspace{-1em}} % mehr Platz vor dem Inhalt

% andere global gemeinsame definitionen
%Includes
\usepackage[ngerman]{babel} %Deutsche Silbentrennung
\usepackage[utf8]{inputenc} %Deutsche Umlaute
\usepackage{float}
\usepackage{graphicx}
\usepackage{minted}
\RequirePackage{csquotes}
\RequirePackage{fontawesome5}

\DeclareGraphicsExtensions{.pdf,.png,.jpg}

\makeatletter
\author{Vorkursteam der Fachschaft Informatik}
\let\Author\@author

% dark mode
\ExplSyntaxOn
\IfDarkModeT{
    \cs_if_exist:NT \setbeamercolor {
        \setbeamercolor*{smallrule}{bg=.}
        \setbeamercolor*{normal~text}{bg=\thepagecolor,fg=.}
        \setbeamercolor*{background~canvas}{parent=normal~text}
        \setbeamercolor*{section~in~toc}{parent=normal~text}
        \setbeamercolor*{subsection~in~toc}{parent=normal~text,fg=.}
        \setbeamercolor*{footline}{parent=normal~text}
        \setbeamercolor{block~title~alerted}{fg=white,bg=white!20!\thepagecolor}
        \setbeamercolor*{block~body}{bg=black!70!gray!98!blue}
        \setbeamercolor*{block~body~alerted}{bg=\thepagecolor}
    }
    \cs_if_exist:NT \setbeamertemplate {
        \setbeamertemplate{subsection~in~toc~shaded}[default][50]
    }
}
\ExplSyntaxOff

% macros
\renewcommand{\arraystretch}{1.2} % Höhe einer Tabellenspalte minimal erhöhen
\newcommand{\N}{{\mathbb N}}
\renewcommand{\code}{\inputminted[]{python}}

\IfDarkModeTF{
    \newmintedfile[pythonfile]{python}{
        fontsize=\small,
        style=native,
        linenos=true,
        numberblanklines=true,
        tabsize=4,
        obeytabs=false,
        breaklines=true,
        autogobble=true,
        encoding="utf8",
        showspaces=false,
        xleftmargin=20pt,
        frame=single,
        framesep=5pt,
    }
    \newmintinline{python}{
        style=native,
        encoding="utf8"
    }
    \newmintinline{kotlin}{
        style=native,
        encoding="utf8"
    }


    \definecolor{codegray}{HTML}{eaf1ff}
    \newminted[bashcode]{awk}{
        escapeinside=||,
        fontsize=\small,
        style=native,
        linenos=true,
        numberblanklines=true,
        tabsize=4,
        obeytabs=false,
        breaklines=true,
        autogobble=true,
        encoding="utf8",
        showspaces=false,
        xleftmargin=20pt,
        frame=single,
        framesep=5pt
    }
}{
    \newmintedfile[pythonfile]{python}{
        fontsize=\small,
        style=friendly,
        linenos=true,
        numberblanklines=true,
        tabsize=4,
        obeytabs=false,
        breaklines=true,
        autogobble=true,
        encoding="utf8",
        showspaces=false,
        xleftmargin=20pt,
        frame=single,
        framesep=5pt,
    }
    \newmintinline{python}{
        style=friendly,
        encoding="utf8"
    }
    \newmintinline{kotlin}{
        style=friendly,
        encoding="utf8"
    }

    \definecolor{codegray}{HTML}{eaf1ff}
    \newminted[bashcode]{awk}{
        escapeinside=||,
        fontsize=\small,
        style=friendly,
        linenos=true,
        numberblanklines=true,
        tabsize=4,
        obeytabs=false,
        breaklines=true,
        autogobble=true,
        encoding="utf8",
        showspaces=false,
        xleftmargin=20pt,
        frame=single,
        framesep=5pt
    }
}

\let\origpythonfile\pythonfile
\renewcommand{\pythonfile}[1]{\pythonfileh{#1}{}}
\newcommand{\pythonfileh}[2]{\origpythonfile[#2]{#1}}

\DeclareDocumentCommand{\kotlinfile}{O{} O{} m}{\inputCode[#1]{minted language=kotlin,#2}{#3}}

\newcommand*{\ditto}{\texttt{\char`\"}}

\newcommand{\shellprefix}{\textcolor{TUDa-3a}{\ttfamily\bfseries \$~}}
\DeclareTCBListing{commandshell}{ O{} O{} }{
    colback=\IfDarkModeTF{black}{black!80},
    colupper=white,
    colframe=TUDa-3a,
    listing only,
    % listing options={style=tcblatex,language=sh},
    listing engine=minted,
    minted style=dracula,
    minted options={
        % linenos=true,
        numbersep=3mm,
        texcl=true,
        autogobble,
        escapeinside=@@,
        breaklines,
        highlightcolor=yellow!50!black,
        #1
    },
    #2,
    % before upper={\textcolor{red}{\small\ttfamily\bfseries root \$> }},
    % every listing line={\textcolor{red}{\small\ttfamily\bfseries root \$> }}
}


%Includes
\usepackage{epstopdf}
\usepackage{wrapfig}
\usepackage{tipa}
\usepackage{tikz}
\usetikzlibrary{calc,shapes,arrows}
%tip: use http://l04.scarfboy.com/coding/phonetic-translation?from=ipa&fromtext=%CB%88pa%C9%AA%CE%B8n%CC%A9&to=tipa
%for converting ipa

\graphicspath{ {./media/} }

\def\shortyear#1{\expandafter\shortyearhelper#1}
\def\shortyearhelper#1#2#3#4{#3#4}

\newcount\NextYear
\NextYear = \year
\advance\NextYear by 1

\newcommand\NextYearShort{\shortyear{\the\NextYear}}

% notes
\usepackage{pgfpages}
\setbeamertemplate{note page}[plain]
%\setbeameroption{show notes on second screen}

% macro for change speaker sign
\newcommand{\changespeaker}{
    \begin{tikzpicture}[line width=.6mm, shorten >= 3pt, shorten <= 3pt]

        \coordinate (c1);
        \coordinate[right of=c1] (c2);

        \draw[rectangle, draw=red!80, fill=red!80, align=center, rounded corners] ($(c1.north west)+(0,-0.3)$) rectangle ($(c2.south east)+(0, 0.3)$) {};
        \draw[->,white] (c1)[bend left] to node[auto] {} (c2);
        \draw[->,white] (c2)[bend left] to node[auto] {} (c1);
    \end{tikzpicture}
}

%Listing-Style pyhon
\title[Programmiervorkurs]{Programmiervorkurs Wintersemester \the\year/\NextYearShort}
\subtitle{{\small der Fachschaft Informatik}}
\logo*{\includegraphics{../globalMedia/bildmarke_ohne_rand}}
\institute{Fachschaft Informatik}
\date{Wintersemester \the\year/\NextYearShort}


% macros
\newcommand{\livecoding}{
    %\ifdefined\StreamSlides
    \begin{frame}
        \frametitle{\insertsectionhead \\  {\small \insertsubsectionhead}}\centering \huge 	\vskip 2cm\textbf{\textcolor{red}{Live-Coding}}
    \end{frame}
    %\fi
}

%\newcommand{\slidehead}{\frametitle{\insertsectionhead \\ {\small \insertsubsectionhead}}\vspace{3mm}}
\newcommand{\slidehead}{\frametitle{\insertsectionhead} \framesubtitle{\insertsubsectionhead}\vspace{3mm}}
\newcommand{\tocslide}{\begin{frame}[t]\frametitle{Inhaltsverzeichnis}\vspace{3mm}{\small\tableofcontents[subsectionstyle=shaded]}\end{frame}}

\newcommand{\nextvid}[2]{
    \ifdefined\StreamSlides
    \else
        \section{Nächstes Video}
        \begin{frame}[t]
            \slidehead
            \begin{block}{Nächstes Video}
                \vspace{0.5cm}
                #1
                \vspace{0.5cm}
            \end{block}
            \ifx\hfuzz#2\hfuzz
                \vspace{2.5cm}
                \else
                {\begin{block}{Bonus Video}
                        \vspace{0.5cm}
                        #2
                        \vspace{0.5cm}
                    \end{block}}
            \fi
            Danke fürs Zuschauen!\\
            Links zu den Folien und Quellen sind in der Videobeschreibung.
        \end{frame}
    \fi
}


\usepackage{verbatim}
\usetikzlibrary{decorations.pathreplacing}
\usetikzlibrary{shapes.misc}


% colors
\definecolor{lightpetrol}{RGB}{\IfDarkModeTF{0,136,119}{0,223,194}}


\usepackage{listings}
\usepackage{tikz}
\usepackage{tabularx,booktabs,multirow,multicol,colortbl}
\aboverulesep = 0mm \belowrulesep = 0mm

\def\streamlink{https://www.youtube.com/d120de/live}
\def\moodlecourselink{https://moodle.informatik.tu-darmstadt.de/course/view.php?id=1636}

\begin{document}

%Deckblatt
\titlegraphic{
    \includegraphics[height=.7\textheight]{bildmarke_ohne_rand.png}
}
\maketitle

\section{Organisatorisches}
\begin{frame}
    \slidehead
    \begin{columns}
        \begin{column}{6cm}
            \begin{figure}
                \centering
                \includegraphics[scale=0.065]{ws14_15}
                \caption{WS 2013/2014}
            \end{figure}
        \end{column}
        \begin{column}{6cm}
            \begin{figure}
                \centering
                \includegraphics[scale=0.21]{ws16_17}
                \caption{WS 2016/2017}
            \end{figure}
        \end{column}
    \end{columns}
\end{frame}

\section{Veranstalter}
\subsection*{Fachschaft Informatik}
\begin{frame}
    \slidehead
    \begin{columns}[b]
        \begin{column}{ .6\textwidth }
            \begin{itemize}[<+->]
                \item \enquote{Die Schüler*innenvertretung der Studierenden}
                \item Wir machen neben Vorkurs \dots
                    \begin{itemize}
                        \item Ophase
                        \item Cloudflight
                        \item Hören gute Musik
                        \item FB20 Parties
                        \item Sommerfest und Nikolausfeier
                    \end{itemize}
                \item Ansprechpartner bei Fragen zum Studium
                \item Gestalten den Studienablauf mit
                    \vspace{.4cm}
                \item {\centering \huge Mehr auf D120.de}
            \end{itemize}
        \end{column}%
        \begin{column}{ .4\textwidth }
            \includegraphics[width=\textwidth]{media/bildmarke_ohne_rand.png}
        \end{column}
    \end{columns}
\end{frame}

\section{Moodle-Kurs}
\begin{frame}[c]
    \slidehead
    \begin{columns}[c]
        \begin{column}{.6\linewidth}
            \centering
            \begin{itemize}
                \item Vorlesungsfolien \& Übungen
                \item Ankündigungen
                \item Fragen, Foren
            \end{itemize}
        \end{column}%
        \begin{column}{.4\linewidth}
            \begin{figure}
                \centering\mbox{}
                \qrcode[height=2.5cm]{\moodlecourselink}
                \caption{\href{\moodlecourselink}{Link zum Moodle-Kurs}}
            \end{figure}
        \end{column}
    \end{columns}
    \vspace{-1em}
    \begin{block}{Lernportal Informatik - Vorkurs}
        {\Huge Gastschlüssel: PVK-2024}

        \vspace{1em}Nähere Informationen zum Zugriff auf die Materialien:\\
        \href{https://d120.de/vorkurs}{https://d120.de/vorkurs}
    \end{block}
\end{frame}


\section{Discord}
\begin{frame}[c]
    \slidehead
    \begin{columns}[c]
        \begin{column}{.6\linewidth}
            \begin{itemize}
                \item Discord-Server der Fachschaft (D120)
                    % \item zu finden im TU Darmstadt Student Hub
                \item dort könnt ihr euch untereinander austauschen
                \item dort könnt ihr auch mit erfahrenen Tutor*innen reden
            \end{itemize}
        \end{column}%
        \begin{column}{.4\linewidth}
            \begin{figure}
                \centering\mbox{}
                % \qrcode[height=2.5cm]{https://discord.gg/JUCns6f7vT}
                % \caption{\href{https://discord.gg/JUCns6f7vT}{Link zum TU-Darmstadt Student Hub}}
                \qrcode[height=2.5cm]{https://discord.gg/GBj9rvszDT}
                \caption{\href{https://discord.gg/GBj9rvszDT}{Link zum D120-Discord-Server}}
            \end{figure}
        \end{column}
    \end{columns}
    \begin{block}{d120 Discord - Vorkurs}
        Für Discord braucht ihr einen Account!

        Ihr könnt Discord hier herunterladen: \url{https://discord.com/}

        % Wie ihr in Discord kommt, erfahrt ihr auf der Vorkurs Seite.
        Lest euch bitte auf dem Server die Willkommensnachricht durch!
    \end{block}
\end{frame}


\section{YouTube}
\begin{frame}[c]
    \slidehead
    \begin{columns}[c]
        \begin{column}{.5\textwidth}
            \begin{itemize}
                \item wir streamen für euch die Vorlesungen auf YouTube live
                \item ihr könnt während dessen fragen auf Discord stellen!
                    %		\item Diskutiert auch gerne parallel in Discord.
                \item Link zum Live-Stream: \url{\streamlink}
            \end{itemize}
        \end{column}%
        \begin{column}{.5\textwidth}
            \begin{figure}
                \centering\mbox{}
                \qrcode[height=3cm]{\streamlink}
                \caption{\href{\streamlink}{Link zum Vorlesungs-Stream}}
            \end{figure}
        \end{column}
    \end{columns}
\end{frame}


\section{Kontakt}
\begin{frame}
    \slidehead
    \begin{itemize}
        \item \textbf{Mail:} \href{mailto:vorkurs@d120.de}{vorkurs@d120.de}
        \item \textbf{Moodle:}  Foren im \href{\moodlecourselink}{Moodle-Kurs}
    \end{itemize}
    \vspace{1cm}
    \begin{defBox}[title=Hinweise]%
        \begin{itemize}
            \item allgemeine Fragen: Schreibt ins Moodle-Forum.
            \item spezielle Fragen an uns: Schreibt uns eine Mail.
            \item Ihr könnt uns natürlich auch ansprechen.
        \end{itemize}
    \end{defBox}
\end{frame}

\section{Ablauf des Vorkurses}
\subsection{Ein üblicher Tag im Vorkurs}
\begin{frame}
    \slidehead
    \begin{itemize}
        \item \textbf{Vormittags Vorlesung (Hybrid) 9:00 - 12:00}
            \begin{itemize}
                \item Vorlesung im Hörsaal S3|11 0012
                \item Übertragung in S2|02 C205, falls zu viele Leute da sind
                \item interaktiv
                \item heißt: Fragen stellen
                \item zusätzlicher Stream auf \url{\streamlink}
            \end{itemize}
            \pause
        \item \textbf{Mittagspause 12:00 - 13:00}
            \begin{itemize}
                \item Pausen sind wichtig
                \item Esst vielleicht ne Kleinigkeit.
                \item Wenn ihr schon eure Athenekarte habt, könnt ihr in die Mensa gehen.
            \end{itemize}
            \pause
        \item \textbf{Nachmittags Übung 13:00 - 18:00}
            \begin{itemize}
                \item Gruppenarbeit im C-Pool
                \item Tutor*innen
                \item offenes Ende
            \end{itemize}
    \end{itemize}
\end{frame}

\subsection{Wochenplan}
\begin{frame}
    \slidehead
    \begin{table}[ht!]
        \small\centering
        \def\vlcolor{\cellcolor{brown!15!\thepagecolor}}
        \def\pausecolor{\cellcolor{accentcolor!15!\thepagecolor}}
        \def\excolor{\cellcolor{TUDa-1b!15!\thepagecolor}}
        \rowcolors{2}{\thepagecolor}{fgcolor!10!\thepagecolor}
        \begin{tabularx}{\textwidth}{l>{\centering\arraybackslash}X>{\centering\arraybackslash}X>{\centering\arraybackslash}X>{\centering\arraybackslash}X>{\centering\arraybackslash}X}
            \toprule
            \fatsf{Zeit} & \fatsf{Montag}& \fatsf{Dienstag}& \fatsf{Mittwoch}& \fatsf{Donnerstag}& \fatsf{Freitag}\\
            \midrule
            \phantom{09:00} &  &  &  &  & \\
            % Vorlesungsfarbe   MO           DI           MI         DO  FR
            \tikzmark{tpos-1} & \vlcolor{} & \vlcolor{} & \vlcolor{} & & \vlcolor{}\\
            \tikzmark{tpos-2} & \vlcolor{} & \vlcolor{} & \vlcolor{} & & \vlcolor{}\\
            \tikzmark{tpos-3} &
            % Vorlesung MO
            \multirow{-3}{*}{\vlcolor{}\begin{minipage}[t][2cm][t]{\textwidth}
                    \tiny \fatsf{Vorlesung}\\
                    Foliensätze:\\
                    \quad - Orga\\
                    \quad - Erste Schritte\\
                    \quad - Operationen\\
                    \quad - Variablen
                \end{minipage}} &
            % Vorlesung DI
            \multirow{-3}{*}{\vlcolor{}\begin{minipage}[t][2cm][t]{\textwidth}
                    \tiny \fatsf{Vorlesung}\\
                    Foliensätze:\\
                    \quad - Schleifen\\
                    \quad - Listen
                \end{minipage}} &
            % Vorlesung MI
            \multirow{-3}{*}{\vlcolor{}\begin{minipage}[t][2cm][t]{\textwidth}
                    \tiny \fatsf{Vorlesung}\\
                    Foliensätze:\\
                    \quad - Funktionen\\
                    \quad - Rekursion\\
                    \quad - Chellange Vorstellung
                \end{minipage}} &
            % Vorlesung DO
            &
            % Vorlesung FR
            \multirow{-3}{*}{\vlcolor{}\begin{minipage}[t][2cm][t]{\textwidth}
                    \tiny \fatsf{Vorlesung}\\
                    Foliensätze:\\
                    \quad - Chellange Lösung\\
                    \quad - Recap\\
                    \quad - Fragestunde
                \end{minipage}}\\
            \tikzmark{tpos-4} &
            %Mittag MO
            \pausecolor{}{\footnotesize Mittagspause} &
            %Mittag DI
            \pausecolor{}{\footnotesize Mittagspause} &
            %Mittag MI
            \pausecolor{}{\footnotesize Mittagspause} &
            %Mittag DO
            &
            %Mittag FR
            \pausecolor{}{\footnotesize Mittagspause}\\
            %Übungsfarbe        MO           DI           MI         DO  FR
            \tikzmark{tpos-5} & \excolor{} & \excolor{} & \excolor{} & & \excolor{}\\
            \tikzmark{tpos-6} & \excolor{} & \excolor{} & \excolor{} & & \excolor{}\\
            \tikzmark{tpos-7} & \excolor{} & \excolor{} & \excolor{} & & \excolor{}\\
            \tikzmark{tpos-8} & \excolor{} & \excolor{} & \excolor{} & & \excolor{}\\
            \tikzmark{tpos-9} &
            % Übung MO
            \multirow{-4}{*}{\excolor{}\begin{minipage}[t][3cm][t]{\textwidth}
                    \tiny \fatsf{Übungsphase}\\
                    \quad - Python installieren\\
                    \quad - Übungsblatt 1
                \end{minipage}} &
            % Übung DI
            \multirow{-4}{*}{\excolor{}\begin{minipage}[t][3cm][t]{\textwidth}
                    \tiny \fatsf{Übungsphase}\\
                    \quad - Übungsblatt 2
                \end{minipage}} &
            % Übung MI
            \multirow{-4}{*}{\excolor{}\begin{minipage}[t][3cm][t]{\textwidth}
                    \tiny \fatsf{Übungsphase}\\
                    \quad - Übungsblatt 3\\
                    \quad - Challenge
                \end{minipage}} &
            % Übung DO
            \multirow{-9}{*}{\textcolor{TUDa-9b}{Feiertag}} &
            % Übung FR
            \multirow{-4}{*}{\excolor{}\begin{minipage}[t][3cm][t]{\textwidth}
                    \tiny \fatsf{Übungsphase}\\
                    \quad - Challenge\\
                    \quad - MRMCD
                \end{minipage}}\\
            \tikzmark{tpos-10} &  &  &  &  & \\
            \bottomrule
        \end{tabularx}

        \begin{tikzpicture}[remember picture,overlay]
            \foreach \x in {1,...,10} {
                \node[anchor=south west] at ([xshift=-1mm,yshift=.5ex]pic cs:tpos-\x) {\padzeroes{\the\numexpr8+\x\relax}:00};
            }
        \end{tikzpicture}
    \end{table}
\end{frame}

\subsection{Voraussetzungen}
\begin{frame}
    \slidehead
    \centering
    \vspace{1.5cm}
    \Huge Keine
\end{frame}

\subsection{Übungsaufgaben}
\begin{frame}
    \slidehead
    \begin{itemize}
        \item Aufgaben mit unterschiedlichen Schwierigkeitsgraden
        \item Tutor*innen
            die euch unterstützen
        \item Helft euch auch gegenseitig, denn so lernt ihr am meisten
        \item Ihr müsst nicht alle Aufgaben lösen
    \end{itemize}
\end{frame}

\section{Ophase}
\begin{frame}
    \slidehead
    \textbf{Nächste Woche} ist die Ophase (07.10. bis 11.10.2023)! \\
    Nehmt auf jeden Fall daran teil. Es lohnt sich für euch!
\end{frame}

\end{document}
